% Generated by scripts/assemble.py; corrections belong in source/corrections.json.
\examhead{184}{长安大学 2024 年数学分析试卷}{173}
% p174-b14; source PDF page 174
\par\Needspace{5\baselineskip}
\blockanchor{p174-b14}
\majorquestion{一}计算题. 每题 5 分, 共 30 分.\par
% p174-b15; source PDF page 174
\blockanchor{p174-b15}
\question{1}求极限 $\lim_{n\to \infty}\left(1 + \frac{1}{n} +\frac{1}{n^2}\right)^n.$\par
% p174-b16; source PDF page 174
\blockanchor{p174-b16}
\question{2}求极限 $\lim_{x\to +\infty}\left(\sqrt[5]{x^5 + x^4} -\sqrt[5]{x^5 - x^4}\right).$\par
% p174-b17; source PDF page 174
\blockanchor{p174-b17}
\question{3}求不定积分 $\int x^{2}e^{-x}\mathrm{d}x.$\par
% p174-b18; source PDF page 174
\blockanchor{p174-b18}
\question{4}$y = \frac{h}{\sqrt{\pi}} e^{-h^2 x^2}$ 在 $x = \pm \delta (\delta >0)$ 处为拐点，求 $h$ 的值 $(h > 0)$\par
% p174-b19; source PDF page 174
\blockanchor{p174-b19}
\question{5}求心形线 $r = a(1 + \cos \theta) (a > 0)$ 所围图形的面积.\par
% p174-b20; source PDF page 174
\blockanchor{p174-b20}
\question{6}求坐标平面与 $x = 2, y = 3, x + y + z = 4$ 所围成的角柱体的体积.\par
% p174-b21; source PDF page 174
\par\Needspace{5\baselineskip}
\blockanchor{p174-b21}
\majorquestion{二}简答题. 每题 10 分, 共 40 分.\par
% p174-b22; source PDF page 174
\blockanchor{p174-b22}
\question{1}证明数列 $\left\{\left(1 + \frac{1}{n + 1}\right)^n\right\}$ 是递增数列, 并求极限 $\lim_{n\to \infty}\left(1 + \frac{1}{n + 1}\right)^n.$\par
% p174-b23; source PDF page 174
\blockanchor{p174-b23}
\question{2}设 $S(x) = \sum_{n=1}^{\infty} ne^{-nx} (x > 0)$, 求 $\int_{\ln 2}^{\ln 3} S(t) \, \mathrm{d}t$.\par
% p174-b24; source PDF page 174
\blockanchor{p174-b24}
\question{3}计算曲面积分 $\iint_{S} (x + y + z) \, \mathrm{d}S$, 其中 $S$ 是上半球面 $x^{2} + y^{2} + z^{2} = a^{2}$, $z \geq 0$.\par
% p174-b25; source PDF page 174
\blockanchor{p174-b25}
\question{4}计算曲线积分 $\int_{L} \frac{x \, \mathrm{d}y - y \, \mathrm{d}x}{x^2 + y^2}$, 其中 $L$ 是 $xOy$ 平面上任一光滑的封闭曲线, 且不过原点, $L$ 的方向为逆时针方向.\par
% p175-b00; source PDF page 175
\par\Needspace{5\baselineskip}
\blockanchor{p175-b00}
\majorquestion{三}(15 分) 设 $f(x)$ 在 $(a, +\infty)$ 上可导, 且 $\lim_{x \to +\infty} f'(x) = 0$, 证明: $\lim_{x \to +\infty} \frac{f(x)}{x} = 0$.\par
% p175-b01; source PDF page 175
\par\Needspace{5\baselineskip}
\par\Needspace{6\baselineskip}
\blockanchor{p175-b01}
\majorquestion{四}(15 分) 设 $f(x)$ 为 $[a, b]$ 上的连续增函数, 记\par
\nopagebreak[4]
% p175-b02; source PDF page 175
\blockanchor{p175-b02}
\displayformula{F (x) = \left\{ \begin{array}{l l} \frac {1}{x - a} \int_ {a} ^ {x} f (t) \mathrm{d} t, & x \in (a, b ]; \\ f (a), & x = a. \end{array} \right.}
% p175-b03; source PDF page 175
\blockanchor{p175-b03}
证明： $F(x)$  为  $[a,b]$  上的单调增函数.\par
% p175-b04; source PDF page 175
\par\Needspace{5\baselineskip}
\blockanchor{p175-b04}
\majorquestion{五}(15 分) 设 $f(x, y) = \begin{cases} \frac{x^2y}{x^2 + y^2}, & x^2 + y^2 \neq 0; \\ 0, & x^2 + y^2 = 0. \end{cases}$\par
% p175-b05; source PDF page 175
\blockanchor{p175-b05}
(1) 讨论 $f(x, y)$ 在 $(0, 0)$ 点的连续性.\par
% p175-b06; source PDF page 175
\blockanchor{p175-b06}
(2) 证明 $f(x, y)$ 在 $(0,0)$ 点偏导数存在.\par
% p175-b07; source PDF page 175
\blockanchor{p175-b07}
(3) 讨论 $f(x, y)$ 在 $(0, 0)$ 点的可微性.\par
% p175-b08; source PDF page 175
\par\Needspace{5\baselineskip}
\blockanchor{p175-b08}
\majorquestion{六}(20 分) 证明含参量广义积分 $\int_{0}^{+\infty} \frac{\sin xy}{y} \, \mathrm{d}y$\par
% p175-b09; source PDF page 175
\blockanchor{p175-b09}
(1) 在  $[\delta, +\infty)$  上一致收敛 ( $\delta > 0$ ).\par
% p175-b10; source PDF page 175
\blockanchor{p175-b10}
(2) 在 $(0, +\infty)$ 上不一致收敛.\par
% p175-b11; source PDF page 175
\par\Needspace{5\baselineskip}
\blockanchor{p175-b11}
\majorquestion{七}(15 分) 设 $f(x)$ 在 $[a, b]$ 上有连续的导函数, 且 $f(a) = 0$.\par
% p175-b12; source PDF page 175
\blockanchor{p175-b12}
(1) 证明: $M^2 \leq (b - a) \int_{a}^{b} [f'(x)]^2 \, \mathrm{d}x$, $M = \max_{a \leq x \leq b} |f(x)|$.\par
% p175-b13; source PDF page 175
\blockanchor{p175-b13}
(2) 证明: $\int_{a}^{b} f^{2}(x) \, \mathrm{d}x \leq \frac{(b - a)^{2}}{2} \int_{a}^{b} [f'(x)]^{2} \, \mathrm{d}x$.\par
